% ============================================================
% Content: Pairs, Photons and Holes
% Torsion conservation as common principle
% behind Cooper pair, exciton and SPDC
% March 2026
%
% Convention of this document:
%   [STANDARD] = established physics with sources
%   [FFGFT]    = interpretation within the FFGFT framework
% ============================================================

\providecommand{\ket}[1]{\left|#1\right\rangle}
\providecommand{\bra}[1]{\left\langle#1\right|}

% ============================================================
\section{Initial Question and Methodology}
\label{sec:initial-question}
% ============================================================

Three phenomena of modern physics show at first glance
a structural similarity:

\begin{itemize}
  \item A photon splits into two photons of lower frequency
    (parametric fluorescence, SPDC).
  \item An excited electron forms a bound pair with a remaining
    hole (exciton).
  \item Two electrons with opposite spin form
    a Cooper pair (superconductivity).
\end{itemize}

In all three cases a bound pair arises from a single
excitation state --- and in all three cases
conservation quantities are exactly satisfied.

\begin{important}[Methodology of this document]
This document strictly separates two levels:

\textbf{[STANDARD]}: Established physics, experimentally secured,
documented with primary sources.
These statements are valid independently of FFGFT.

\textbf{[FFGFT]}: Interpretation within the FFGFT framework.
These statements are FFGFT-specific and
should be read as such.

The aim is to rigorously establish the structural analogy between the three systems --- and then to specify precisely where the FFGFT interpretation goes beyond standard physics.
\end{important}

% ============================================================
\section{System 1: Parametric Photon Splitting (SPDC)}
\label{sec:spdc}
% ============================================================

\subsection{[STANDARD] Established Physics of SPDC}

\textbf{Spontaneous Parametric Down-Conversion} (SPDC) is a
nonlinear optical process: a photon with frequency $\omega_0$
is converted in a nonlinear crystal into two photons
of lower frequency~\cite{Burnham1970}:
\begin{equation}
  \omega_0 \;\longrightarrow\; \omega_1 + \omega_2,
  \quad \text{with } \omega_1 + \omega_2 = \omega_0.
\end{equation}

Exactly two conservation laws hold:

\textbf{Energy conservation:}
\begin{equation}
  \hbar\omega_0 \;=\; \hbar\omega_1 + \hbar\omega_2.
\end{equation}

\textbf{Momentum conservation (phase matching):}
\begin{equation}
  \mathbf{k}_0 \;=\; \mathbf{k}_1 + \mathbf{k}_2.
\end{equation}

For type-II SPDC the resulting photons have
\textbf{opposite polarisation} and are \textbf{entangled}.
The total angular momentum of the system is conserved~\cite{Kwiat1995}.

\begin{keyresult}[STANDARD: What SPDC directly shows]
\begin{enumerate}
  \item A higher-energy state splits into two
    lower-energy states.
  \item Energy and momentum are exactly conserved.
  \item The opposite polarisation of the photons
    is not an additional assumption --- it follows from
    angular momentum conservation.
  \item The resulting photons are \textbf{real photons}
    --- not quasi-particles, not computational aids.
    \emph{[FFGFT caveat: see below --- photons in FFGFT
    are not point particles but extended wave packets.]}
\end{enumerate}
\end{keyresult}

\subsection{[FFGFT] Interpretation of SPDC}

\begin{foundation}[FFGFT: SPDC as torsion splitting --- photons as wave packets]
In FFGFT every photon carries angular momentum as a
topological property of the electromagnetic field torsion
(Doc.~174, Section~3).

\textbf{Important FFGFT basic position: photons are wave packets.}

In standard theory (QED), the photon is a quantised
excitation quantum of the electromagnetic field --- mathematically
a Fock state $\ket{1_\mathbf{k}}$.
It is often treated as a point particle.

In FFGFT the photon is \textbf{not a point particle},
but an \textbf{extended wave packet}:
a spatially bounded, coherent torsion wave of the
electromagnetic field with defined frequency $\omega$,
wave vector $\mathbf{k}$ and polarisation.

\textbf{Why this matters for SPDC:}
The splitting $\omega_0 \to \omega_1 + \omega_2$ is
in FFGFT not the decomposition of a point particle,
but the \textbf{splitting of a wave packet}
into two wave packets of lower frequency in the nonlinear crystal.
This is physically more direct --- light \emph{is} a wave,
and the splitting is a wave process.

\begin{center}
\begin{tabular}{p{5cm}p{6cm}}
\toprule
\textbf{Standard theory (QED)} & \textbf{FFGFT} \\
\midrule
Photon = Fock state $\ket{1_\mathbf{k}}$ &
  Photon = extended torsion wave packet \\
Point particle picture for many calculations &
  No point particle; always wave packet \\
Splitting = annihilation + creation of quanta &
  Splitting = wave process in nonlinear medium \\
Entanglement = quantum correlation of Fock states &
  Entanglement = topological torsion correlation \\
\bottomrule
\end{tabular}
\end{center}

\textbf{Consequence for the description ``real photon'':}
``Real'' means in FFGFT: no quasi-particle approximation,
no many-body reformulation --- but no point particle.
A SPDC photon is a real extended wave packet
of the torsion field with clearly defined properties
($\omega$, $\mathbf{k}$, polarisation, spatial extent).

\textbf{Conserved quantity in FFGFT:}
The topological torsion charge is additive and conserved.
This is no new claim --- it is the FFGFT formulation
of the known angular momentum conservation, now formulated
for wave packets instead of point particles.
\end{foundation}

% ============================================================
\section{System 2: The Exciton --- Electron and Hole}
\label{sec:exciton}
% ============================================================

\subsection{[STANDARD] What a Hole Really Is}

Here lies the conceptual core.

\begin{important}[STANDARD: The hole is not a real particle]
A \textbf{hole} in the valence band of a semiconductor is
\textbf{not an independent fundamental particle} ---
it is a \textbf{quasi-particle} that describes the collective
reaction of the many-body system to a missing electron~\cite{Kittel2004,AshcroftMermin1976}.

Precisely formulated (standard textbook definition):
\begin{itemize}
  \item The valence band is completely filled in the ground state.
  \item When an electron is missing (through absorption of a photon),
    the many-body system behaves \emph{as if}
    a positively charged particle with positive effective mass
    were present.
  \item This fictitious particle is called a ``hole''.
  \item It has: positive charge $+e$, effective mass $m_h^*$,
    and a spin opposite to that of the missing electron.
\end{itemize}

The hole is mathematically equivalent to the description
in electron language --- but it is a construct of the
many-body approximation, not a real antiparticle
(not a positron!).
\end{important}

The real antiparticle of the electron is the \textbf{positron} ---
which arises in pair production ($\gamma \to e^- + e^+$)
from real high-energy photons. The hole in the semiconductor
is structurally different: it exists only in the context
of the valence band of the material.

\subsection{[STANDARD] The Exciton}

The \textbf{exciton} is a bound state consisting of an
excited electron in the conduction band and a hole
in the valence band~\cite{Wannier1937,Frenkel1931}:
\begin{equation}
  \text{Exciton} = \text{Electron (conduction band)} + \text{Hole (valence band)}.
\end{equation}

The binding arises through the Coulomb attraction between the
negatively charged electron and the positively charged hole.

\textbf{Two main types:}
\begin{itemize}
  \item \textbf{Wannier-Mott exciton}: weakly bound
    (binding energy 1--100\,meV), spatially extended
    (radius $\gg a_\text{lattice}$). Typical for semiconductors
    (GaAs, CdSe, Si). Stable only at low temperatures.
  \item \textbf{Frenkel exciton}: strongly bound
    (binding energy 0.1--1\,eV), spatially compact
    (radius $\sim a_\text{lattice}$). Typical for organic
    semiconductors, biological systems. Stable even at room temperature.
\end{itemize}

\subsection{[FFGFT] Interpretation of the Exciton}

\begin{foundation}[FFGFT: Exciton as complementary torsion pair]
In FFGFT the electron in the conduction band carries
a definite torsion configuration (spin, $\xi$-level $N$).
The hole in the valence band is described as
anti-torsion: the absence of an electron with
opposite quantum numbers.

The exciton is in FFGFT a \textbf{torsion-anti-torsion pair}:
the total torsion is reduced (Spin $S_\text{total}$,
Momentum $\mathbf{p}_\text{total}$ are conserved),
analogously to SPDC.

\textbf{Critical limitation:}
The hole is a quasi-particle in standard theory ---
no fundamental particle. The FFGFT description as
``anti-torsion'' is an interpretation that goes beyond
standard physics and has not been independently proven.
\end{foundation}

% ============================================================
\section{System 3: The Cooper Pair}
\label{sec:cooper-pair}
% ============================================================

\subsection{[STANDARD] BCS Theory and the Cooper Pair}

In the BCS theory of superconductivity~\cite{BCS1957,Cooper1956},
two electrons with opposite spin and momentum
form a bound pair:
\begin{equation}
  \text{Cooper pair} = \text{Electron}(\mathbf{k},\uparrow) +
  \text{Electron}(-\mathbf{k},\downarrow).
\end{equation}

The binding is not direct (two negative charges repel each other)
but is mediated via phonons (lattice vibrations):
\begin{itemize}
  \item Electron 1 polarises the crystal lattice.
  \item This polarisation attracts Electron 2 with a delay.
  \item The result is an indirect, weak attractive interaction.
\end{itemize}

Binding energy: $\Delta E \approx 1$--$10$\,meV
(conventional superconductors). This is much smaller than
$k_B T$ at room temperature ($\approx 26$\,meV).

\begin{center}
\resizebox{\textwidth}{!}{%
\begin{tabular}{p{3cm}p{3cm}p{3cm}p{3cm}}
\toprule
\textbf{Type} & \textbf{$T_c$} & \textbf{Binding energy} & \textbf{Mechanism} \\
\midrule
Conventional & 1--40\,K & 1--10\,meV & Phonon-mediated \\
High-$T_c$ & up to 138\,K & $\sim 50$--100\,meV & Not fully understood \\
\bottomrule
\end{tabular}
}
\end{center}

\subsection{[FFGFT] Interpretation of the Cooper Pair}

\begin{foundation}[FFGFT: Cooper pair as opposite torsion configurations]
In FFGFT the two electrons of the Cooper pair have
opposite torsion configurations:
$(\mathbf{k}, \uparrow)$ carries torsion in one direction,
$(-\mathbf{k}, \downarrow)$ carries torsion in the opposite direction.

The pair thus has:
\begin{itemize}
  \item Total torsion zero (spin singlet)
  \item Total momentum zero ($\mathbf{k} + (-\mathbf{k}) = 0$)
\end{itemize}

This is a \textbf{topologically neutral state} ---
geometrically more stable than two free, separate torsion configurations.

\textbf{Why low temperature is necessary (FFGFT interpretation):}
The binding does not arise directly from the torsion geometry,
but is mediated via phonons (lattice vibrations).
Phonons are collective vibration modes of the crystal lattice ---
they lie at a different $\xi$-level than the electron spin system.
The phonon-mediated coupling is therefore a
cross-level coupling between $\xi$-levels
(Doc.~175, section on $\xi$-hierarchy),
whose strength $\propto \xipar^{\Delta N}$ with $\Delta N = 1$.
This explains the weakness of the binding energy
compared to direct spin-spin coupling.

\textbf{Critical limitation:}
This interpretation is consistent but not proven.
Standard BCS theory explains the same situation
completely without FFGFT concepts.
\end{foundation}

% ============================================================
\section{Comparison of the Three Systems}
\label{sec:comparison}
% ============================================================

\subsection{[STANDARD] Structural Commonalities}

\begin{center}
\resizebox{\textwidth}{!}{%
\begin{tabular}{p{2.5cm}p{3cm}p{3cm}p{3cm}p{2.5cm}}
\toprule
\textbf{System} & \textbf{Constituents} & \textbf{Binding energy} &
\textbf{Temperature} & \textbf{Conserved} \\
\midrule
\textbf{SPDC} &
  Photon~1 + Photon~2 &
  --- (no binding) &
  Room temperature &
  $E$, $\mathbf{p}$, $J_z$ \\[6pt]
\textbf{Exciton} &
  Electron (conduction band) + Hole (valence band, quasi-particle) &
  $10$--$100$\,meV (Wannier) to 1\,eV (Frenkel) &
  4\,K--300\,K depending on type &
  $E$, $\mathbf{p}$, $S$, charge \\[6pt]
\textbf{Cooper pair} &
  Electron~1 ($\mathbf{k},\uparrow$) + Electron~2 ($-\mathbf{k},\downarrow$) &
  $1$--$10$\,meV (conventional) &
  $< T_c$ (1--40\,K) &
  $E$, $\mathbf{p}=0$, $S=0$ \\
\bottomrule
\end{tabular}
}
\end{center}

\textbf{Structural commonality [STANDARD]:}
In all three cases, two objects with opposite quantum numbers
(spin, momentum, polarisation) arise from an excited state,
so that conservation quantities are exactly satisfied.

\textbf{Critical difference [STANDARD]:}
\begin{itemize}
  \item \textbf{SPDC:} two real photons, no binding energy,
    no temperature limit.
  \item \textbf{Exciton:} real electron + quasi-particle (hole),
    Coulomb binding, temperature-sensitive (Wannier-Mott: $T \ll T_\text{bind}/k_B$).
  \item \textbf{Cooper pair:} two real electrons, phonon-mediated
    binding, very temperature-sensitive ($T < T_c$).
\end{itemize}

The hole in the exciton is \textbf{not} analogous to the second photon
in SPDC --- it is a quasi-particle without independent
fundamental existence outside the many-body system.

\subsection{[STANDARD] The User's Question: ``Hole = Two Electrons?''}

This is a precise and important question.

The answer is: \textbf{technically yes, conceptually no}.

\textbf{Technically yes:}
The exciton description ``electron + hole'' is mathematically
equivalent to a description exclusively in electrons.
In many-body theory the hole is defined as:
\begin{equation}
  \ket{\text{hole}} \;\equiv\; \ket{\text{Fermi sea without 1 electron at }
  \mathbf{k}, \downarrow}.
\end{equation}
The hole \emph{is} the missing electron in disguise.

\textbf{Conceptually no:}
The hole description is not merely a renaming ---
it is an efficient reorganisation of the many-body states
that describes many phenomena (conductivity, magnetotransport,
optical transitions) much more simply than pure
electron language.
The hole \emph{behaves} like a particle --- it has
an effective mass, a spin, a momentum.
It is a useful construct with real physical content,
even if it is not a fundamental particle.

\begin{keypoint}[Hole: useful construct, not a fundamental particle]
The hole in the semiconductor is not:
\begin{itemize}
  \item A real antiparticle (not a positron)
  \item A new fundamental structure of nature
  \item Independent of the electron description
\end{itemize}

The hole is:
\begin{itemize}
  \item A quasi-particle: an emergent excitation of the
    many-body system
  \item Mathematically: the absence of an electron in the
    valence band, rewritten as a positively charged particle
  \item Physically meaningful: it correctly describes real,
    measurable phenomena
\end{itemize}

The \textbf{Cooper pair}, by contrast, consists of
\textbf{two real electrons} --- no hole,
no quasi-particle. That is the fundamental difference.
\end{keypoint}

% ============================================================
\section{[FFGFT] The Common Principle: Torsion Conservation}
\label{sec:torsion-conservation}
% ============================================================

\begin{foundation}[FFGFT: Conservation of total torsion as a unified principle]
In FFGFT the connecting principle behind all three systems is
the \textbf{conservation of total torsion}:

\begin{center}
\resizebox{\textwidth}{!}{%
\begin{tabular}{p{2.5cm}p{4cm}p{4cm}p{3cm}}
\toprule
\textbf{System} & \textbf{Before the process} &
\textbf{After the process} & \textbf{Conserved} \\
\midrule
SPDC &
  Torsion $\omega_0$, angular momentum $J_z$ &
  Torsion $\omega_1$ (+) + Anti-torsion $\omega_2$ (--) &
  Torsion charge: $J_z = J_1 + J_2$ \\[6pt]
Exciton &
  Valence band torsion (full band) &
  Conduction band torsion (electron) + valence band anti-torsion (hole) &
  Torsion charge: $S_\text{total}$, $\mathbf{p}_\text{total}$ \\[6pt]
Cooper pair &
  Fermi sea torsion & Pair torsion ($\mathbf{k},\uparrow$) + ($-\mathbf{k},\downarrow$) &
  Torsion charge: $S = 0$, $\mathbf{p} = 0$ \\
\bottomrule
\end{tabular}
}
\end{center}

\textbf{The common FFGFT language:}
In all three cases a pair of torsion and anti-torsion arises
from a torsion configuration, such that the
total torsion charge is conserved.

This is in FFGFT no coincidence --- it is the mirror image
of the known conservation laws (energy, momentum, angular momentum)
in the language of the torsion field.

\textbf{What FFGFT does not claim:}
FFGFT does \emph{not} claim that exciton hole and
Cooper electron and SPDC photon are the same objects ---
they are physically different.
FFGFT claims that their \emph{structural similarity}
has a common geometric origin:
the conservation of the topological torsion charge.
\end{foundation}

% ============================================================
\section{Consistency with the Temperature Argument from Doc.~173}
\label{sec:temperature}
% ============================================================

The three systems differ dramatically in their
temperature stability. This is [STANDARD] and directly explicable:

\begin{center}
\resizebox{\textwidth}{!}{%
\begin{tabular}{p{2.8cm}p{2.5cm}p{3cm}p{4cm}}
\toprule
\textbf{System} & \textbf{Binding energy} &
\textbf{Stability temperature} & \textbf{Cause} \\
\midrule
Photon (SPDC) &
  none (free photons) &
  Room temperature &
  No thermal dissipation \\[6pt]
Frenkel exciton &
  0.1--1\,eV &
  Room temperature possible &
  $\Delta E \gg k_BT(300\,\text{K}) = 26$\,meV \\[6pt]
Wannier exciton (CdSe) &
  $\sim 15$--20\,meV &
  $< 50$\,K &
  $\Delta E \sim k_BT$ at room temperature \\[6pt]
Cooper pair (conventional) &
  1--10\,meV &
  $T < T_c = 1$--40\,K &
  $\Delta E \ll k_BT(300\,\text{K})$ \\
\bottomrule
\end{tabular}
}
\end{center}

\begin{keypoint}[Consistency with Condition 2 from Doc.~173]
The temperature limits follow directly from Condition~2
(Doc.~173): $\Delta E > k_BT$.
The weaker the binding energy, the lower
the temperature must be --- this is no FFGFT argument,
this is standard physics.

The Cooper pair needs low temperature not because of
its electron nature, but because of its
\textbf{very weak binding energy},
which follows from the indirect phonon mechanism.
\end{keypoint}

% ============================================================
\section{Summary}
\label{sec:summary-179}
% ============================================================

\begin{conclusionBox}[Conclusion: What is certain and what is FFGFT interpretation]

\textbf{[STANDARD] Established statements:}

\textbf{1. The hole is a quasi-particle, not a real antiparticle.}\\
It is the mathematical reformulation of a missing electron
in the valence band. ``Exciton consists of two electrons'' is
technically correct --- but the hole description is more efficient
and physically meaningful.

\textbf{2. The Cooper pair consists of two real electrons.}\\
Opposite spin and momentum. Phonon-mediated binding.
Low temperature is necessary because of the weak binding energy
($\Delta E \ll k_BT_\text{room temperature}$).

\textbf{3. SPDC produces two real photons --- in FFGFT: two wave packets.}\\
Energy, momentum and angular momentum are exactly conserved.
No binding energy --- the photons are free.
In FFGFT photons are not point particles but
extended torsion wave packets --- the splitting
is a wave process, not the decay of a point particle.

\textbf{4. The structural similarity is real:}\\
In all three cases two objects with
opposite quantum numbers arise from an initial state,
so that conservation quantities are exactly satisfied.

\medskip

\textbf{[FFGFT] Interpretation:}

\textbf{5. Torsion conservation as a unified principle.}\\
The known conservation laws (energy, momentum, angular momentum)
correspond in FFGFT to the conservation of the topological
torsion charge. The structural similarity of the three systems
thus has a common geometric origin ---
without the three systems being the same.

\textbf{6. The hole as anti-torsion.}\\
Consistent with the FFGFT description of the exciton
from Doc.~173 --- but going beyond standard physics
and not independently proven.
\end{conclusionBox}

% ============================================================

\begin{thebibliography}{99}

% --- SPDC ---
\bibitem{Burnham1970}
Burnham, D.\,C.\ \& Weinberg, D.\,L.
\textit{Observation of simultaneity in parametric production of
optical photon pairs.}
Phys.\ Rev.\ Lett.\ \textbf{25}, 84--87 (1970).

\bibitem{Kwiat1995}
Kwiat, P.\,G.\ et al.
\textit{New high-intensity source of polarization-entangled photon pairs.}
Phys.\ Rev.\ Lett.\ \textbf{75}, 4337--4341 (1995).

% --- Exciton ---
\bibitem{Frenkel1931}
Frenkel, J.
\textit{On the transformation of light into heat in solids. I \& II.}
Phys.\ Rev.\ \textbf{37}, 17--44 and 1276--1294 (1931).

\bibitem{Wannier1937}
Wannier, G.\,H.
\textit{The structure of electronic excitation levels in
insulating crystals.}
Phys.\ Rev.\ \textbf{52}, 191--197 (1937).

% --- Semiconductor fundamentals (hole as quasi-particle) ---
\bibitem{Kittel2004}
Kittel, C.
\textit{Introduction to Solid State Physics.}
8th ed.\ Wiley (2004). Chapter~8: Semiconductor Crystals.

\bibitem{AshcroftMermin1976}
Ashcroft, N.\,W.\ \& Mermin, N.\,D.
\textit{Solid State Physics.}
Holt, Rinehart and Winston (1976). Chapter~12: The Semiclassical
Model of Electron Dynamics; Chapter~28: Homogeneous Semiconductors.

% --- Cooper pair and BCS ---
\bibitem{Cooper1956}
Cooper, L.\,N.
\textit{Bound electron pairs in a degenerate Fermi gas.}
Phys.\ Rev.\ \textbf{104}, 1189--1190 (1956).

\bibitem{BCS1957}
Bardeen, J., Cooper, L.\,N.\ \& Schrieffer, J.\,R.
\textit{Theory of superconductivity.}
Phys.\ Rev.\ \textbf{108}, 1175--1204 (1957).

% --- Exciton in quantum dots ---
\bibitem{Woggon1997}
Woggon, U.
\textit{Optical Properties of Semiconductor Quantum Dots.}
Springer Tracts in Modern Physics, Vol.\ 136 (1997).

\end{thebibliography}
